\hypertarget{group__macro__for__factories}{}\section{Macros to simplify insertion in factories}
\label{group__macro__for__factories}\index{Macros to simplify insertion in factories@{Macros to simplify insertion in factories}}
The five macros

S\+M\+Spp\+\_\+insert\+\_\+in\+\_\+factory\+\_\+h

S\+M\+Spp\+\_\+insert\+\_\+in\+\_\+factory\+\_\+cpp\+\_\+0( $<$ class name $>$ )

S\+M\+Spp\+\_\+insert\+\_\+in\+\_\+factory\+\_\+cpp\+\_\+1( $<$ class name $>$ )

S\+M\+Spp\+\_\+insert\+\_\+in\+\_\+factory\+\_\+cpp\+\_\+0\+\_\+t( $<$ class name $>$ )

S\+M\+Spp\+\_\+insert\+\_\+in\+\_\+factory\+\_\+cpp\+\_\+1\+\_\+t( $<$ class name $>$ )

can be used to quickly ensure that $<$ class name $>$ is inserted in the corresponding factory. As the name says, they have to be put respectively in the private part of $<$ class name $>$ declaration in the .h, and in the .cpp of $<$ class name $>$ (if any, {\itshape exactly one} .cpp otherwise). The \char`\"{}\+\_\+k\char`\"{} versions of the \+\_\+cpp macro refer to the fact that the constructor of the base class has k parameters. The further \char`\"{}\+\_\+t\char`\"{} versions need be used if the class is template, since then a template specialization is needed and this requires adding \char`\"{}template$<$$>$\char`\"{} at proper places. 